\documentclass[10pt,a4paper]{amsart}
\usepackage[margin=19mm]{geometry}
\usepackage{amsmath,amssymb,mathtools}
\usepackage[scr=rsfs,cal=euler]{mathalfa}
\usepackage{xfrac}
\usepackage[unicode,hidelinks,bookmarks=false]{hyperref}
\newcommand{\ie}{i.e.\ }
\newcommand{\cf}{cf.\ }
\newcommand{\resp}{respectively\ }
\newcommand{\Subsection}{\subsection}
\providecommand{\nobreakdash}{\nobreak\hbox{-}}
\setlength{\emergencystretch}{2em}
\multlinegap0pt
\usepackage{amssymb}
\usepackage{xcolor}
\newtheorem{lemma}{Lemma}
\newcommand{\R}{{\mathbb R}}
\newcommand{\la}{\lambda}
\newcommand{\supp}{\mathrm{supp }}
\title{Lossless Strichartz and spectral projection estimates on unbounded manifolds}
\author{Xiaoqi Huang, Christopher D. Sogge, Zhongkai Tao, Zhexing Zhang}
\date{Source: arXiv:2504.07238v3 (CC BY 4.0)}
\begin{document}
\maketitle

\noindent\emph{Source and licence.} Lossless Strichartz and spectral projection estimates on unbounded manifolds. Version arXiv:2504.07238v3 (CC BY 4.0), distributed by arXiv under the Creative Commons Attribution 4.0 International licence, http://creativecommons.org/licenses/by/4.0/, as recorded in the arXiv licence metadata for this exact version and shown on the abstract page https://arxiv.org/abs/2504.07238v3. Full licence notice: \texttt{licenses/arxiv-2504.07238-CC-BY-4.0.txt}.\par\smallskip

\noindent\textit{Editorial scope.} Counted unit: the article's lemma co (source0.tex lines 775--795). The article's complete original proof of it is delivered with it (source0.tex lines 988--1114, one \texttt{proof} environment of 127 lines, which the article places away from the statement). No mathematics is written, repaired or re-derived: the statement and the proof are the article's own lines, in the article's own notation, and no retained line is substituted or re-flowed.  No further source-local context is needed: every cross-reference of every retained line resolves inside the two delivered spans.  Nothing was omitted from any retained span, so the package carries no omission record.  No retained line contains a citation, so the wrapper carries no bibliography and no bibliography entry.  No retained line contains a figure environment, an image-inclusion command or drawing code, so no image is delivered and the package declares no asset dependency.  Editorial formatting only, in addition: an AMS \texttt{amsart} preamble that restates the article's own declarations and macros used by the retained lines, namely the environments \texttt{lemma} on separate counters, together with the standard packages those lines need.  The retained environments print in this document as Lemma 1 in order of appearance, whereas the article prints the same items with the numbering of its own full document.  The complete native source of the article as distributed in the arXiv e-print for this exact version is bundled unchanged as \texttt{sources/176275/source0.tex}.
\begin{lemma}\label{co}
    Let $\chi\in C_0^\infty(M)$ satisfy $\chi=1$ in a neighborhood of $\pi(K)$. Let $\beta\in C_0^{\infty}((1/2,2))$ be a Littlewood-Paley bump function satisfying $\sum_{k=-\infty}^\infty \beta(s/2^k)=1$,
$s>0$. Suppose $\tilde \chi\in C_0^\infty(M)$ supported away from the trapped set $\pi(K)$ and $\tilde\chi=1$ in a neighborhood of $\supp \nabla\chi$,  and let $\tilde \beta \in C^\infty_0((1/4,4))$ with $\tilde \beta=1$ in $(1/3,3)$, we have for $\la\gg 1$,
\begin{equation}\label{commuteerror}
    \chi \cdot \beta(\sqrt{-\Delta_g}/\la)f = \tilde\beta(\sqrt{-\Delta_g}/\la) \chi \cdot \beta(\sqrt{-\Delta_g}/\la)  f
+R_1f,
\end{equation}
as well as  
\begin{equation}\label{commuteerror1}
    [\Delta_g,\chi] \cdot \beta(\sqrt{-\Delta_g}/\la)f= \tilde\beta(\sqrt{-\Delta_g}/\la)  [\Delta_g,\chi] \cdot \beta(\sqrt{-\Delta_g}/\la)\tilde\chi \,\tilde\beta(\sqrt{-\Delta_g}/\la)  f
+R_2f,
\end{equation}
where for $i=1,2$,
\begin{equation}\label{r}
    \|(-\Delta_g)^{\alpha}R_if\|_{L^2(M)}\le C_{\alpha,N}\la^{-N}\|f\|_{L^2(M)}  \,\,\forall \alpha>0, N=1,2,3\cdots.
\end{equation}
Moreover,
\begin{equation}\label{commuteerror2}
\| [\Delta_g,\chi] \cdot \beta(\sqrt{-\Delta_g}/\la)f\|_{L^2}\lesssim \la\|f\|_{L^2}.
\end{equation}
\end{lemma}

\begin{proof}[Proof of Lemma~\ref{co}]


We first prove \eqref{commuteerror2}. It suffices to show for $\la_1\ge 8\la$
\begin{equation}\label{c1}
\| \beta(\sqrt{-\Delta_g}/\la_1)[\Delta_g,\chi] \cdot \beta(\sqrt{-\Delta_g}/\la)f\|_{L^2}\lesssim_N \la_1^{-N}\|f\|_{L^2}, 
\end{equation}
as well as 
\begin{equation}\label{c2}
\| \beta_0(\sqrt{-\Delta_g}/\la)[\Delta_g,\chi] \cdot \beta(\sqrt{-\Delta_g}/\la)f\|_{L^2}\lesssim \la\|f\|_{L^2},
\end{equation}
for any $\beta_0\in C_0^\infty(-10,10)$. 

To prove \eqref{c1}, 
 we can extended $\beta$ to be an even function by adding the negative part if necessary. Let $P=\sqrt{-\Delta_g}$, and let $\delta_0$ be fixed and smaller than the injectivity radius of $M$. 
Fix an even function $\rho\in C_0^\infty$   satisfying $\rho(t)= 1$, $|t|\le \delta_0/4$ and $\rho(t)=0 $, $|t|\ge \delta_0/2$ such that
\begin{equation}\label{c3}
\begin{aligned}
     \beta( P/\la)=&(2\pi)^{-1}\int_\R\la\hat \beta(\la t)\cos t P dt \\
     =&(2\pi)^{-1}\int\rho(t)\la\hat \beta(\la t)\cos t P dt+ (2\pi)^{-1}\int (1-\rho(t)) \la\hat \beta(\la t)\cos t P dt\\
     =& B_\la(P)+C_\la(P).
\end{aligned}
\end{equation}
It is not hard to check that the symbol of the operator $C_\la$ is $O_N((1+|\tau|+|\la|)^{-N})$, which implies that 
  \begin{equation}\label{cestimatee}
      \|\Delta_g^{N_0} C_\la(P)\|_{L^2\to L^2}\lesssim_{N_0,N}\la^{-N},
  \end{equation}
 for any fixed $N_0, N$.


  
  
   To handle the term in \eqref{c1} involving $C_{\la}(P)$, note that 
\begin{equation*}
    \begin{split}
         \|\beta(P/\la_1) [\Delta_g, \chi] C_\la(P) f\|_{L^2}
         &\lesssim \la_1^{-2N}\|\beta(P/\la_1)\Delta_g^N [\Delta_g, \chi] C_\la(P) f\|_{L^2}\\
         &\lesssim \la_1^{-2N}\|B_{\la_1}(P)\Delta_g^N [\Delta_g, \chi] C_\la(P) f\|_{L^2}\\
         &\quad+\la_1^{-2N}\|C_{\la_1}(P)\Delta_g^N [\Delta_g, \chi] C_\la(P) f\|_{L^2}.
    \end{split}
\end{equation*}
By \eqref{cestimatee}, the second term on the right side is bounded by
\begin{equation}\label{c4}
\begin{split}
&\la_1^{-2N}\|C_{\la_1}(P)\Delta_g^N [\Delta_g, \chi] C_\la(P) f\|_{L^2} 
\\ &\leq
    \la_1^{-2N}\|C_{\la_1}(P)\Delta_g^{N+1}  \chi C_\la(P) f\|_{L^2} 
    +  \la_1^{-2N}\|C_{\la_1}(P)\Delta_g^N \chi\,\Delta_g C_\la(P) f\|_{L^2}\\
    &\lesssim_N \la_1^{-N}\| f\|_{L^2}.
\end{split}
\end{equation}
For the first term, we have
\begin{equation}\label{c5}
    \begin{split}
    &\la_1^{-2N}\|B_{\la_1}(P)\Delta_g^N [\Delta_g, \chi] C_\la(P) f\|_{L^2} \\
    &\lesssim \la_1^{-2N}\sum\| B_{\la_1}(P)[\Delta_g,[\Delta_g,[\cdots,\chi]\cdots]\Delta_g^{k} C_{\lambda} (P)f\|_{L^2} \\
    &\lesssim \la_1^{-N+1}\sum_{k\le N}\|\Delta_g^k C_{\lambda} (P)f\|_{L^2}\\
    &\lesssim \la_1^{-N+1}\|f\|_{L^2}.
    \end{split}
\end{equation}
Here we used the fact that 
\begin{equation}\label{c6}
    \| B_{\la_1}(P)[\Delta_g,[\Delta_g,[\cdots,\chi]\cdots]\|_{L^2\to L^2}\lesssim \la_1^{N+1}
\end{equation}
if the number of brackets is bounded by $N+1$. By using the parametrix construction for $\cos tP$, 
 one can argue as in the compact manifold case to show that the local operator $B_{\la_1}$ is a 0-order pseudodifferential operator with principal symbol $\beta(p(x,\xi)/\la_1)$,
with $p(x,\xi)$ here being the principal symbol of $P$. Thus it is not hard to show  $B_{\la_1}(P)[\Delta_g,[\Delta_g,[\cdots,\chi]\cdots]$ is a pseudodifferential operator of order $N+1$ with symbol supported in the region $p(x,\xi)\approx \la_1$, which implies \eqref{c6}.

   To handle the term in \eqref{c1} involving $B_{\la}(P)$, note that
\begin{equation*}
\begin{split}
     \| \beta(P/\lambda_1)[\Delta_g, \chi]  B_{\la}(P)) f\|_{L^2}
\leq 
         \| B_{\la_1}(P)[\Delta_g, \chi]  B_{\la}(P)) f\|_{L^2}  +    \|C_{\la_1}(P)[\Delta_g, \chi]  B_{\la}(P)) f\|_{L^2}. 
\end{split}
\end{equation*}
The second term can be estimated as in \eqref{c4}.
Hence it suffices to show
\begin{equation}\label{c7}
    \|B_{\la_1}(P)[\Delta_g, \chi] B_{\la}(P)\|_{L^2\to L^2}\lesssim \la_1^{-N}.
\end{equation}
This follows from integration by parts after writing out the composition of the corresponding pseudodifferential operators explicitly.

To prove \eqref{c2},  write $  \beta_0( P/\la)
     = B^0_\la(P)+C^0_\la(P)$ as in \eqref{c3}. Repeating the argument above, it suffices to show that
\begin{equation}\label{c8}
    \|B^0_{\la}(P)[\Delta_g, \chi] B_{\la}(P)\|_{L^2\to L^2}\lesssim \la.
\end{equation}
This follows from the fact that the operator above is a pseudodifferential operator of order 1 with symbol supported in the region $p(x,\xi)\lesssim\la$.

To prove \eqref{commuteerror1}, it suffices to show 
\begin{equation}\label{c9}
    [\Delta_g,\chi] \cdot \beta(P/\la)f= \tilde\beta(P/\la)  [\Delta_g,\chi] \cdot \beta(P/\la)f
+R_1f,
\end{equation}
as well as 
\begin{equation}\label{c10}
   \tilde\beta(P/\la)  [\Delta_g,\chi] \cdot \beta(P/\la)f
   = \tilde\beta(\sqrt{-\Delta_g}/\la)  [\Delta_g,\chi] \cdot \beta(P/\la)\tilde\chi \,\tilde\beta(P/\la) f
+R_2f,
\end{equation}
where the operators $R_1$ and $R_2$ satisfy \eqref{r}.

As above, to prove \eqref{c9}, it suffices to show for $\la_1\ge 8\la$
\begin{equation}\label{c1a}
\| \beta(P/\la_1)(1-\tilde\beta)(P/\la)[\Delta_g,\chi] \cdot \beta(P/\la)f\|_{L^2}\lesssim_N \la_1^{-N}\|f\|_{L^2}, 
\end{equation}
as well as 
\begin{equation}\label{c2a}
\| \beta_0(P/\la)(1-\tilde\beta)(P/\la)[\Delta_g,\chi] \cdot \beta(P/\la)f\|_{L^2}\lesssim_N \la^{-N}\|f\|_{L^2},
\end{equation}
for any $\beta_0\in C_0^\infty(-10,10)$. 
Both \eqref{c1a} and \eqref{c2a} follow from the same argument used in the proof of \eqref{c1}.

To prove \eqref{c10}, note that by using \eqref{commuteerror2} and duality, it is not hard to show 
$$\|(-\Delta_g)^\alpha\tilde\beta(P/\la)  [\Delta_g,\chi]\|_{L^2\to L^2}\lesssim \la^{2\alpha}\|\tilde\beta(P/\la)  [\Delta_g,\chi]\|_{L^2\to L^2} \lesssim \la^{1+2\alpha}$$
Hence \eqref{c10} is a consequence of 
\begin{equation}\label{c12}
\chi_0 \cdot \beta(P/\la)f
   =\chi_0 \cdot \beta(P/\la)\tilde\chi \,\tilde\beta(P/\la) f
+Rf,
\end{equation}
if $\chi_0\in C_0^\infty$ is chosen so that $\chi_0=1$ on $\supp\, \nabla\chi$, and $\tilde \chi=1$ in a $\delta_0$-neighborhood of the support of $\chi_0$, with $\delta_0$ defined as in \eqref{c3}. Here we may assume $\chi_0$ is supported in a sufficiently small neighborhood of $\supp\, \nabla\chi$, and 
 $\delta_0$ is chosen small enough so that $\tilde \chi$ is supported away from the trapped set $\pi(K)$.

The proof of \eqref{c12}, as well as the remaining estimate \eqref{commuteerror} in Lemma~\ref{co}, follows from a simpler but analogous argument as above, together with the fact that the operator $B_\la(P)$ in \eqref{c3} is local due to finite propagation speed of the wave propagator $\cos(tP)$. We omit the details for simplicity.
\end{proof}

\end{document}
